% from ega0-19-fr.tex
\oldpage[0\textsubscript{IV-1}]{106}

mais $\varprojlim_j((B/\mathfrak m^i)\otimes_A(C/\mathfrak n'{}^j))$ est le
séparé complété de $(B/\mathfrak m^i)\otimes_A C$, et comme $C$ est complet,
ce n'est autre que $(B/\mathfrak m^i)\otimes_A C$ lui-même,
$B/\mathfrak m^i$ étant un $A$-module de type fini puisque
$\mathfrak m/\mathfrak m^i$ est un $(B/\mathfrak m)$-module de type fini
($\mathbf 0_{\mathrm I}$, 7.3.6). On a donc
\[
  E=\varprojlim_i((B/\mathfrak m^i)\otimes_A C).
\]
Pour tout entier $h>0$, $\mathfrak m^hE$, étant un idéal de $E$, est fermé
dans $E$ ($\mathbf 0_{\mathrm I}$, 7.3.5), donc complet, et d'autre part il est
évidemment dense dans
\[
  \varprojlim_i
  (\operatorname{Im}(\mathfrak m^h/\mathfrak m^{h+i})\otimes_A C),
\]
donc égal à cette dernière limite projective. En outre, tous les systèmes
projectifs considérés sont définis par des homomorphismes \emph{surjectifs}~;
il résulte donc de ($\mathbf 0_{\mathrm{III}}$, 13.2.2) que
$E/\mathfrak m^hE$ est isomorphe à
\[
  (B/\mathfrak m^h)\otimes_A C
  =(B\otimes_A C)/\operatorname{Im}(\mathfrak m^h\otimes_A C).
\]
En particulier, comme
$\operatorname{Im}(\mathfrak m\otimes_A C)\subset\mathfrak r$, cela montre
que $\mathfrak mE$ est contenu dans $\mathfrak rE$, donc dans le radical de
$E$. Enfin, l'hypothèse que $C$ est un $A$-module plat entraîne que
$(B/\mathfrak m^h)\otimes_A C=E/\mathfrak m^hE$ est un
$(B/\mathfrak m^h)$-module plat pour tout $h>0$~: comme $B$ et $E$ sont
noethériens et que $\mathfrak mE$ est contenu dans le radical de $E$, il
résulte de ($\mathbf 0_{\mathrm{III}}$, 10.2.2) que $E$ est un $B$-module
plat.
\end{proof}

\begin{lemma}[19.7.1.3]
\label{0.19.7.1.3-fr}
Soient $A$ un anneau local noethérien, $\mathfrak m$ son idéal maximal, $k$ son
corps résiduel, $B_0$ une $k$-algèbre~; on suppose que $B_0$ est un anneau
local noethérien, complet et régulier. Alors il existe une $A$-algèbre
topologique $B$ qui est un anneau local noethérien complet, un $A$-module plat,
et tel que $B_0$ soit $k$-isomorphe à
$B\otimes_A k=B/\mathfrak mB$.
\end{lemma}

\begin{proof}
Comme $\widehat A$ est plat sur $A$ et a même corps résiduel, on peut se borner
au cas où $A$ est \emph{complet}.

Soit $K$ le corps résiduel de $B_0$, et distinguons deux cas~:

\medskip
\noindent I) $K$ est une extension \emph{séparable} de $k$. En vertu de
(19.6.4), $B_0$ est $k$-isomorphe à un anneau de séries formelles
$K[[T_1,\ldots,T_n]]$. Lorsque $B_0=K$, le lemme a déjà été démontré
($\mathbf 0_{\mathrm{III}}$, 10.3.1)~; soit $C$ un anneau local noethérien
complet qui est un $A$-module plat et tel que $C\otimes_A k$ soit isomorphe à
$K$. Pour $n\geq1$, il suffit de prendre (avec la notation précédente)
$B=C[[T_1,\ldots,T_n]]$~; on sait en effet (Bourbaki,
\emph{Alg. comm.}, chap.~III, §~3, n\textsuperscript{o}~4, cor.~3 du th.~1)
que $B$ est un $C$-module plat, donc aussi un $A$-module plat, et d'autre part,
il est immédiat que $C[[T_1,\ldots,T_n]]\otimes_A k$ est isomorphe à
$(C/\mathfrak mC)[[T_1,\ldots,T_n]]=B_0$.

\medskip
\noindent II) $K$ est de \emph{caractéristique} $p>0$, et par suite il en est
de même de $k$. Notons $P$ le corps premier $\mathbf F_p$, et $W(P)$ l'anneau
local complet des nombres $p$-adiques $\mathbf Z_p$, qui est un anneau de
valuation discrète (donc \emph{régulier}), et a $P$ pour corps résiduel.
Montrons d'abord qu'il existe un homomorphisme continu d'anneaux $W(P)\to A$
faisant donc de $A$ une $W(P)$-\emph{algèbre topologique}. En effet, si
$j:\mathbf Z\to A$ est l'homomorphisme canonique, on a
$j(p\mathbf Z)\subset\mathfrak m$ par hypothèse, d'où
$j^{-1}(\mathfrak m)=p\mathbf Z$, et par suite $j$ se factorise en
$\mathbf Z\to\mathbf Z_{p\mathbf Z}\xrightarrow{f}A$, où $f$ est un
homomorphisme local, donc continu, qui (puisque $A$ est complet) se prolonge
par continuité en l'homomorphisme $W(P)\to A$ cherché.

Comme $k$ est extension séparable de $P$, le cas I) montre qu'il y a un
homomorphisme local $W(P)\to W(k)$, où $W(k)$ est un anneau local noethérien
complet et un $W(P)$-module plat, tel que $W(k)\otimes_{W(P)}P$ soit isomorphe
à $k$. D'ailleurs, comme l'uniformisante $p$ de $W(P)$ est un élément
$W(k)$-régulier par platitude ($\mathbf 0_{\mathrm I}$, 6.3.4) et comme

